\section{Threads \& Processes}

\begin{definition}

                A process is a program executing inside an OS (operating-system). Processes share CPU but have seperate
                memory space.
            
\end{definition}

\begin{definition}

                The \textbf{context} of a process is the set of utilities it uses and or operates on. That is:
                \begin{itemize}
    \item \textbf{Hardware context: }Instruction counter, Registers, Execution State, Flags ...
    \item \textbf{Memory context: }Virtual Memory, Heap-Space, Stack ...
    \item \textbf{OS-level context:} Process ID, open files ...
\end{itemize}

\end{definition}

            The lifecylce of a process is managed by the OS and can be described by the following states:
            \begin{itemize}
    \item \textbf{CREATED: }The process is loaded from the disc
    \item \textbf{WAITING: }The process is waiting in the main-memory for execution.
    \item \textbf{RUNNING: }The process is executed
    \item \textbf{BLOCKED: }The process is waiting for access to I/O devices
    \item \textbf{TERMINATED: }The process has finished its execution
\end{itemize}

\begin{definition}

\textbf{Context Switch:}\\

                The process of saving the context of a process and loading the context of another process. This happens
                for example when a process recives the RUNNING state and another process gets BLOCKED.
            
\end{definition}

            All processes are managed by the OS, which schedules, synchronizes, starts and terminates processes, allows
            comminication and manages ressources. The OS uses so called Process control blocks (PCB).

            \begin{definition}

                A \textbf{process-control-block (PSB)} is a table-entry storing relevant process information such as
                Program counter, Registers, Adress Space, Open files ...
            
\end{definition}

\subsection{Parallelism}

\begin{definition}

\textbf{Parallelism } is when multiple processes execute simultaneously on diffrent CPU-Cores.
            
\end{definition}

\begin{definition}

\textbf{Concurrency } is when multiple processes execute interleaved on a single CPU-Core.
            
\end{definition}

            Observe that parallelism implies concurrency, but concurrency does not imply parallelism.

            \subsection{Multithreading}

\begin{definition}

                A thread is an independent sequence of exection running in the same OS-process.
            
\end{definition}

\begin{remark}

                While processes can only communicate through main-memory and have diffrent memory spaces, threads share
                the same adress space. This leads to faster and more efficient context-switches as adressspace and PCB
                loading is not necessary.
            
\end{remark}

            There are diffrent levels of threads:
            \begin{itemize}
    \item \textbf{User-Level Threads: } Threads managed by the application using a thread library, ex.
                    Java-Threads managed by JVM
    \item \textbf{Kernel-Level Threads: }Threads managed by to OS, which can schedule them on diffrent cores.
                
    \item \textbf{CPU-Level Threads:} Threads managed by the CPU, which can schedule them on the physical cores.
                    Ex. A CPU might have 4 physical cores, which apear to the OS as 8 logical cores.
\end{itemize}

            There are two ways to create a thread in java.

            \begin{codeexample}

                class MyThread extends Thread { ...
                public void run() {
                ...
                }
                }
                MyThread thread = new MyThread();
                thread.start(); // calls MyThread.run()
            
\end{codeexample}

            or using a runnable, which is recomended to do:

            \begin{codeexample}

                class MyRunnable implements Runnable { ...
                public void run() {
                ...
                }
                }
                MyRunnable runnable = new MyRunnable();
                Thread thread = new Thread(runnable);
                thread.start(); // calls MyRunnable.run()
            
\end{codeexample}

            In its lifecylce a thread can have the follwoing states:
            \begin{itemize}
    \item \textbf{NEW: }The thread has been created but not yet started.
    \item \textbf{RUNNABLE: }The thread is ready to run and waiting for the OS to schedule it.
    \item \textbf{RUNNING: }The thread is currently executing.
    \item \textbf{BLOCKED: }The thread is waiting for access to I/O devices
    \item \textbf{WAITING: }The thread is waiting to be notified
    \item \textbf{TERMINATED: }The thread has finished its execution
\end{itemize}

\begin{remark}

                Every Thread has an \textbf{ID, Name, Priority and State} which might be usfull for certain tasks.
            
\end{remark}

            To wait for a thread to finish, we can use the .join() method on the thread we want to wait on. This lets
            the main-thread sleep until the joined thread terminates and wakes it up again.

            \begin{codeexample}

                Thread t = new MyThread();
                t.start(); //start t
                t.join(); //sleep until t is finished
                continueWork(); // this is only called after t finishes
            
\end{codeexample}

\begin{remark}

                Using the .join() method might not always be the best option, as it produceses a significant overhead.
                If the execution time of a thread is short, it might be better to let the main-thread busy wait in a
                while loop.
            
\end{remark}

\begin{important}

                If a worker-thread throws an exception the main-thread will be woken up, but it will not know of the
                error,
                which might lead to false results. To make sure we don't produce false results we have to set an
                UncaughtExeptionHandler using a try-catch block when joining threads. The same holds for threads that
                run indefinetly we need to make sure to have a timer task that interupts threads or manually interupt
                them.
            
\end{important}

\begin{codeexample}

                Thread t = new MyThread();
                Thread timer = new MyTimerThread();
                t.start();
                timer.start(t, 1000);
                try {
                t.join();
                } catch (UncaughtExeptionHandler) {
                System.out.println("Error in thread t");
                }
                continueWork();

                MyTimerThread extends Thread {
                run(Thread thread, int time) {
                timer.sleep(time);
                thread.interupt();
                }
                }
            
\end{codeexample}

